\section{Registros de Decisiones Arquitectónicas (ADR)}

\begin{figure}[htbp]
\centering
\includegraphics[width=\textwidth]{"../../architecture/SDK/Hexagonal architecture class diagram.png"}
\caption{Diagrama de la arquitectura hexagonal}
\label{fig:hexagonal}
\end{figure}

\subsection{ADR-01: Adopción de Arquitectura Hexagonal para el núcleo del SDK}
\textbf{ASR asociado:} ASR-01, Extensibilidad ante nuevas pasarelas de pago

\textbf{Problema.} El proyecto nace de un problema muy concreto: cuatro pasarelas de pago colombianas (Wompi, Rapyd, Mercado Pago y Kushki) exponen contratos completamente distintos entre sí, con diferencias en autenticación, formato de solicitud, nomenclatura de estados y mecanismos de firma. Si el SDK acoplara su lógica de negocio directamente a las particularidades de cada proveedor, cualquier cambio en una API externa, o la necesidad de incorporar una quinta pasarela en el futuro, obligaría a modificar el núcleo del sistema, incrementando el riesgo de introducir errores en funcionalidades que no tenían relación con el cambio.

\textbf{Solución.} Se organizó el SDK siguiendo el estilo de Arquitectura Hexagonal, también conocido como Puertos y Adaptadores, separando el código en tres capas independientes: dominio, aplicación e infraestructura. El dominio concentra las reglas de negocio y no conoce ninguna pasarela en particular; la comunicación con el exterior ocurre exclusivamente a través del puerto PaymentGatewayPort, definido en la capa de aplicación, mientras que cada proveedor se traduce mediante un adaptador independiente (WompiAdapter, RapydAdapter, MercadoPagoAdapter y KushkiAdapter) que implementa dicho puerto siguiendo el patrón GoF Adapter. Esta organización puede verse en el Diagrama de Componentes del SDK y en el Diagrama de Clases de la Arquitectura Hexagonal (Figura 16), donde se aprecia que ningún componente del dominio referencia directamente un tipo o estructura propia de Wompi, Rapyd, Mercado Pago o Kushki.

\textbf{Consecuencias.} Como consecuencia positiva, agregar una nueva pasarela se convierte en una tarea acotada a la creación de un nuevo adaptador, sin tocar el dominio, la fachada ni los adaptadores existentes, lo que además reduce directamente el acoplamiento (CBO) medido sobre el código cliente del SDK. Como consecuencia negativa, esta separación introduce una mayor cantidad de archivos y una curva de entendimiento inicial más alta para quien no esté familiarizado con el estilo hexagonal, en comparación con una arquitectura en capas tradicional donde el código estaría más concentrado aunque más acoplado.

\subsection{ADR-02: Adopción de Arquitectura Hexagonal para el núcleo del SDK}
\textbf{ASR asociado:} ASR-02, Cambio de pasarela activa sin modificar código cliente.

\textbf{Problema.} Un comercio que integra el SDK necesita poder cambiar de pasarela de pago, por ejemplo pasar de Wompi a Rapyd, sin que ese cambio implique reescribir su lógica de negocio ni recompilar partes del sistema que no tienen relación directa con la pasarela. Si la selección del proveedor quedara dispersa entre distintos componentes, o si el desarrollador tuviera que instanciar manualmente el adaptador correspondiente, el cambio de pasarela dejaría de ser una operación de configuración y se convertiría en un cambio de código con riesgo de introducir errores.

\textbf{Solución.} Se introdujo el SDK Configurator como punto único de gestión de la configuración, responsable de mantener la pasarela activa, las credenciales por proveedor, el modo de entorno y los parámetros del mecanismo de reintento. A partir de esa configuración, el Gateway Factory, implementado bajo el patrón GoF Factory, recibe el valor del enum Gateway y retorna la instancia del adaptador correspondiente sin que el resto del sistema conozca ese proceso de instanciación. El Payment Facade, como único punto de entrada del SDK, consulta al Configurator y delega en el Factory antes de ejecutar cualquier operación, tal como se describe en el Diagrama de Componentes del SDK.

\textbf{Consecuencias.} La consecuencia positiva más directa es que cambiar de pasarela se reduce a modificar un parámetro de inicialización del SDK, sin que el desarrollador que lo consume necesite conocer ni modificar la lógica interna de selección de adaptadores. La consecuencia negativa es que el Gateway Factory sigue siendo el componente que debe actualizarse cada vez que se incorpora una nueva pasarela al enum Gateway, por lo que, aunque el acoplamiento se reduce drásticamente para el código cliente, no desaparece por completo dentro del propio SDK.

\subsection{ADR-03: Adopción de Arquitectura Hexagonal para el núcleo del SDK}
\textbf{ASR asociado:} ASR-03, Normalización uniforme de respuestas entre pasarelas

\textbf{Problema.} Cada pasarela devuelve el resultado de una transacción con su propio vocabulario: Wompi usa APPROVED, Rapyd usa el valor CLO acompañado de un campo booleano paid Mercado Pago responde en minúsculas con approved, y Kushki utiliza el término APPROVAL en lugar de APPROVED para un pago exitoso. Si el desarrollador que integra el SDK tuviera que interpretar directamente estos valores nativos, su código quedaría acoplado a los formatos de cada proveedor y expuesto a errores sutiles, como el caso documentado de Kushki, donde una comparación literal contra APPROVED clasificaría incorrectamente todos los pagos exitosos.

\textbf{Solución.} Se definió el Response Normalizer como servicio de aplicación encargado de traducir el estado nativo devuelto por cualquier adaptador hacia el enum canónico TransactionStatus, con los valores APPROVED, DECLINED, PENDING, EXPIRED, VOIDED y ERROR. Este servicio construye además la entidad de dominio Transaction, que expone métodos de negocio como isApproved() o isFinal(), de modo que el desarrollador razone sobre el resultado de un pago sin conocer el formato original de la pasarela. El valor nativo se conserva en el campo rawStatus para efectos de auditoría, siguiendo la tabla de equivalencias documentada en el lenguaje ubicuo del proyecto. El flujo completo, desde la respuesta nativa hasta la entidad normalizada, está representado en el Diagrama de Secuencia de Creación de Pago.

\textbf{Consecuencias.} Entre las consecuencias positivas está que el código cliente del SDK nunca necesita referenciar los tipos de respuesta de ninguna pasarela, lo que reduce su acoplamiento y hace que la lógica de negocio del comercio sea independiente del proveedor activo. Entre las consecuencias negativas está que la tabla de mapeo de estados debe mantenerse y ampliarse manualmente cada vez que se agrega una pasarela nueva o que un proveedor existente modifica su vocabulario de estados, concentrando en este único componente un riesgo de error que antes estaba distribuido entre varios adaptadores.

\subsection{ADR-04: Verificación criptográfica de firmas para webhooks entrantes}
\textbf{Problema.} Las notificaciones asíncronas de las pasarelas llegan como solicitudes HTTP hacia un endpoint del comercio, y nada impide, en principio, que un tercero malintencionado lance una solicitud fabricada que simule un pago aprobado si no existe un mecanismo de verificación. El reto adicional es que cada pasarela firma sus webhooks de forma distinta, Wompi mediante SHA-256 sobre una cadena de propiedades más un secreto, Rapyd mediante HMAC-SHA256 sobre una cadena que incluye la URL pública configurada para el webhook, una marca de tiempo y las credenciales del comercio, codificando el resultado en Base64, y Mercado Pago mediante HMAC-SHA256 sobre encabezados y cuerpo.

\textbf{Solución.} Se definió el Webhook Verifier como servicio de aplicación con un único método verify(payload, headers, secret, gateway), que internamente delega en la implementación de firma correspondiente a cada proveedor. Este componente se apoya en el mismo puerto de dominio para mantener el aislamiento respecto a las particularidades de cada pasarela, de forma consistente con la Arquitectura Hexagonal adoptada en el ADR-01, y su interacción con el resto del sistema se documenta en el Diagrama de Secuencia de Conciliación de Webhooks. Solo después de que la firma es validada exitosamente, el resultado se delega hacia la actualización del estado de la entidad Transaction.

\textbf{Consecuencias.} La consecuencia positiva es que el SDK ofrece una única función de verificación independiente del proveedor, evitando que el comercio implemente su propia lógica criptográfica y reduciendo la superficie de error ante ataques de suplantación de eventos. La consecuencia negativa es que la robustez de esta verificación depende directamente de la documentación pública de cada pasarela.

\subsection{ADR-05: Reintentos con backoff exponencial ante fallas transitorias}
\textbf{ASR asociado:} ASR-05, Recuperación ante fallas transitorias de red

\textbf{Problema.} Las integraciones con servicios externos están expuestas a fallos que no reflejan un problema de negocio, como timeouts de red o respuestas HTTP 5xx del lado de la pasarela. Si el SDK propagara inmediatamente cualquiera de estos errores hacia el desarrollador, una simple congestión momentánea en la infraestructura del proveedor terminaría reportándose como un fallo definitivo, obligando al comercio a implementar su propia lógica de reintentos de forma dispersa y potencialmente inconsistente entre pasarelas.

\textbf{Solución.} Se incorporó el Retry Handler como servicio de aplicación que intercepta los errores antes de que lleguen al Error Handler, distinguiendo entre errores transitorios candidatos a reintento, como timeouts o códigos 5xx, y errores definitivos que no se recuperarán con intentos adicionales, como un rechazo de pago o credenciales inválidas. La estrategia adoptada es un backoff configurable, con tres reintentos por defecto e intervalos de uno, dos y cuatro segundos. Si se agota el número máximo de intentos, el error se entrega al Error Handler con el código MAX\_RETRIES\_EXCEEDED.

\textbf{Consecuencias.} La consecuencia positiva es que el SDK absorbe automáticamente una clase completa de fallas de infraestructura que, de otro modo, degradarían la experiencia de integración sin ninguna relación con la lógica de negocio del comercio. La consecuencia negativa es que introduce latencia adicional en escenarios donde el reintento efectivamente no tendrá éxito, y que la clasificación entre error transitorio y definitivo depende de reglas fijas que podrían no cubrir todos los casos particulares de cada pasarela conforme estas evolucionen.

\subsection{ADR-06: Motor de escenarios controlados en la API de Simulación}
\textbf{ASR asociado:} ASR-06, Reproducción fiel de escenarios de prueba en el simulador

\textbf{Problema.} Validar el comportamiento del SDK frente a los distintos resultados posibles de un pago (aprobado, rechazado, fondos insuficientes, timeout o caída del servidor) requeriría, en un escenario real, provocar esas condiciones directamente sobre las pasarelas de producción o sandbox, lo cual no siempre es posible ni deseable durante la fase de evaluación académica del proyecto. Sin un mecanismo controlado, el equipo no tendría forma confiable ni repetible de ejercitar cada rama de la lógica del SDK, en particular el Response Normalizer y el Retry Handler.

\textbf{Solución.} Se construyó la API de Simulación sobre Fastify y TypeScript, aprovechando su soporte nativo de validación mediante JSON Schema para definir con precisión la forma de las solicitudes y respuestas de cada escenario. El flujo se organizó como una tubería de procesamiento secuencial: el HTTP Router identifica la pasarela a partir de la URL e intercepta los encabezados \texttt{x-simulate-scenario} y \texttt{x-simulate-delay}, el Scenario Execution Engine decide cómo proceder según el escenario solicitado, y el Gateway Mock Factory, junto con sus fábricas internas por proveedor, construye un payload que replica con exactitud la estructura nativa de Wompi, Rapyd, Mercado Pago o Kushki según corresponda. Este diseño está documentado en el Diagrama de Componentes de la API de Simulación, y su efecto observable se ilustra en los Diagramas de Secuencia de Pago Exitoso, Pago Rechazado y Error de Red.

\textbf{Consecuencias.} La consecuencia positiva es que se puede probar de forma determinista y repetible cada camino relevante del SDK sin depender de la disponibilidad ni del comportamiento real de las cuatro pasarelas externas, lo cual resulta particularmente valioso dado el contexto de evaluación académica del proyecto. La consecuencia negativa es que la validez de las pruebas depende por completo de la fidelidad de los mocks frente a las respuestas reales de cada proveedor, de modo que si una pasarela modifica su formato de respuesta y esa actualización no se refleja en el Gateway Mock Factory, las pruebas seguirían pasando aunque el SDK ya no fuera compatible con el comportamiento real de dicha pasarela.

